\section{Limitaciones}
\label{sec:limitaciones}

Se enumeran por escrito, en el informe y en el repositorio, porque conocerlas
es parte del trabajo y porque la alternativa, dejarlas implícitas, convertiría
el proyecto en algo que no es.

\subsection{Del sistema}

\begin{description}[leftmargin=0.6cm,style=nextline]
  \item[Un solo nodo de base de datos.] Sin alta disponibilidad ni
    replicación. Es un punto único de fallo y el techo de escritura es el de
    una instancia.
  \item[La limitación de peticiones cuenta en memoria del proceso.] Con cuatro
    workers, la ráfaga efectiva del inicio de sesión es veinte y no cinco.
    Resolverlo exige un contador compartido.
  \item[Tokens sin revocación inmediata.] Un token robado es válido hasta que
    caduca. Se mitiga con una vida de quince minutos; la solución completa es
    una lista de revocación o sesiones en base de datos.
  \item[El token del navegador vive en \codigo{sessionStorage}.] Se borra al
    cerrar la pestaña, pero cualquier ataque de tipo \emph{cross-site
    scripting} podría leerlo. La respuesta correcta en producción es una
    cookie \codigo{httpOnly} con protección contra falsificación de petición
    entre sitios.
  \item[Una cuenta por usuario.] El modelo admite varias, pero el MVP impone
    una por unicidad.
  \item[Docker Compose no es un despliegue de producción.] No hay
    orquestación, ni secretos gestionados, ni cifrado en reposo.
  \item[Sin observabilidad continua.] Hay registro estructurado y correlación,
    pero no métricas exportadas ni trazas distribuidas.
\end{description}

\subsection{De las mediciones}

\begin{description}[leftmargin=0.6cm,style=nextline]
  \item[Cliente y servidor en la misma máquina.] El generador de carga, los
    cuatro workers y PostgreSQL compiten por los mismos dieciséis hilos de
    CPU. Los números absolutos serían distintos con la carga generada desde
    fuera; las comparaciones entre estrategias, que es para lo que se
    hicieron, sí son válidas porque todas sufren la misma interferencia.
  \item[Omisión coordinada.] Locust es un generador de bucle cerrado: cuando
    el servidor se ralentiza deja de emitir peticiones, de modo que las
    latencias altas quedan subrepresentadas. Las cifras de cola son por tanto
    optimistas.
  \item[Un minuto por carrera.] Suficiente para el régimen estacionario,
    insuficiente para ver efectos lentos como el crecimiento de las tablas o
    el comportamiento del recolector de versiones.
  \item[Sin medición de recursos.] No se registró uso de CPU ni de memoria por
    proceso. Queda como pendiente, no como dato estimado.
  \item[Una sola repetición por combinación.] No hay intervalos de confianza.
    Las diferencias que se discuten son lo bastante grandes como para no
    depender del ruido, pero las pequeñas no deberían interpretarse.
  \item[La carga ambiente varía entre carreras.] Una misma configuración
    medida en dos momentos distintos dio 280,6 y 154,6 peticiones por segundo
    en el escenario de lectura. Por eso las comparaciones del
    apartado~\ref{subsec:optimizacion} se hicieron con las dos carreras
    seguidas.
  \item[No es una prueba de capacidad.] No se buscó el punto de saturación ni
    la carga máxima sostenible.
\end{description}

\subsection{Del alcance}

PayCore no implementa conciliación con entidades externas, ni reversiones
contables, ni límites regulatorios, ni prevención de blanqueo, ni
notificaciones, ni recuperación de contraseña, ni verificación de identidad.
Nada de eso está a medias: directamente no existe, y decirlo es más útil que
insinuar lo contrario.
